\documentclass[a4paper]{article}
\usepackage[dvinames]{xcolor}
\usepackage[top=2in,left=1.5in,bottom=2.8in,right=1.25in]{geometry}
\usepackage{graphicx}
\usepackage{ctable}
\usepackage[colorlinks=false,
            pdfborder={0 0 0},]{hyperref}
\usepackage[absolute]{textpos}
\usepackage[oldstylenums]{kpfonts}
\usepackage[T1]{fontenc}
\usepackage{fancyhdr}
\pagestyle{fancy}
\parindent0pt
\setlength{\parskip}{0.8\baselineskip}
\definecolor{SRblue}{RGB}{13,31,97}
\definecolor{SRgrey}{RGB}{223,219,215}
\urlstyle{sf}
\fancyhead{}
\fancyfoot{}
\setlength{\footskip}{3em}
\renewcommand{\headrulewidth}{0pt}
\renewcommand{\footrulewidth}{1pt}
\fancyfoot[C]{\footnotesize\color{SRblue}\sffamily
\url{info@studentrobotics.org} \textbullet\ \url{http://studentrobotics.org/}\\ \medskip \Huge
\makebox[0.8\textwidth]{\includegraphics[height=1em]{moto.png}\hspace{-3em}\vbox{\vfill{}\small Student Robotics is made possible\\ by our generous sponsors\vfill{}}\hspace{-3em}\includegraphics[height=1em]{ecs.png}}}
\begin{document}
\begin{textblock*}{2in}[0.3066,0.39](1.5in,1.33in)
    \includegraphics[width=1.7in]{sr-logo.pdf}
\end{textblock*}
\begin{textblock*}{6.375in}(0.65in,1.0in)
    \sffamily
    \begin{flushright}
    \sffamily \Huge \color{SRblue}
    Student Robotics\\ \large
    \smallskip
    info@studentrobotics.org\\
    https://www.studentrobotics.org\\
    \end{flushright}
\end{textblock*}
\color{black}
\vspace*{0.4in}

\section*{Welcome to Student Robotics}

Thank you for registering for Student Robotics! This document is a guide intended to
be read by teachers or other adults who are submitting a team to the competition. It
covers how the competition is structured, what goes on throughout the Student Robotics
year, what we provide, and what we require from you.


\subsection*{Who's who}

The people who run Student Robotics (often referred to as ``the blueshirts'')
are a collection of engineers, scientists and students who share a passion for
promoting engineering and education for young people. All volunteers, we work
together to help bring the experience of a substantial and creative engineering
project to those considering studying STEM subjects or working in the industry.

The competitors are any 16-18 year old pupils who are interested in learning
about robotics, electronics and programming. We don't require that they all
come from the same school or college, or are even in education, only that
they're interested in learning and building a robot. Together, in a group of at
least three competitors, they form a team who'll be building a robot to compete
at our competition.

You should be a teacher, parent, or other adult who is authorised to take
responsibility for the competitors. You'll be coordinating with us to get the
team to events and the competition, and supporting the team with any local
logistics (for example access to tools or work space). You'll also need to
enforce some of the requirements on competitors mentioned later in this
document.

\subsection*{Keeping in contact}

The primary method that we use to provide support to teams is through mentors
and via the website forum, see below. If you have any non-technical questions
or enquiries that are unanswered or not suitable for those support channels,
please contact our organising committee at \url{info@studentrobotics.org}. If
you have a particularly urgent enquiry, please contact Jeremy Morse on
07XXXXXXXXX.

\section*{What Student Robotics provides}

To help the team build a robot for our competition, Student Robotics provides
an electronics kit that makes controlling devices and sensors simple. The kit,
in combination with a vision system and a Python programming environment,
makes it easy to build a self-controlled robot. The exact pieces of equipment
in the kit change from year to year, but details can be found in the ``kit''
section of the documentation on our
website.\footnote{\url{https://www.studentrobotics.org/docs/kit/}}
We also provide detailed documentation of what the equipment is capable of,
and how to go about programming it.

Occasionally parts of the electronics kit break, sometimes due to competitor
error, and sometimes not. If you think part of the kit is broken,
please report it on the forum. We will diagnose what's wrong and
send a replacement if it's broken -- you'll need to post us the broken piece
back. Accidents do happen, and we don't mind replacing some kit, however if
the team frequently damages equipment we may stop replacing parts until we
work out why.

A list of what piecese of equipment you will need to return at the end of the
competition will be in the competition rule book published at Kickstart.

Our electronics kit is on loan between Kickstart and the competition. All of the
kit must be returned at the competition, unless you arrange to extend the loan
for a longer period (please see the rule book once it is published for details).
Unfortunately we are not in a position to let anyone permanently keep the kit.

\subsection*{Online resources}

As well as the documentation for our equipment, we also provide an online
forum\footnote{\url{https://www.studentrobotics.org/forum}} for competitors,
team leaders and blueshirts to discuss the competition, ask technical questions,
and receive support.

We also provide an online development
environment\footnote{\url{https://www.studentrobotics.org/ide}} for competitors
to write robot code, and export it to their robot. This development environment
(``The IDE'') has a code checker to catch errors early, and keeps all the
teams robot code in one place to avoid code getting lost. We also occasionally
ship software updates for the robot via the IDE.

The website is also used for competitors to return media consent forms, and
get their ticket to the competition. We require each competitor to have their
own account on the website -- accounts are created by you, the team leader,
through the user management
page.\footnote{\url{https://www.studentrobotics.org/userman/}} Documentation
for using the user management page is XXX I need to upload this.

\subsection*{Events}

There are two primary events that we organise in the Student Robotics year --
Kickstart introduces Student Robotics to the competitors, and reveals what this
years game will be. The kit is distributed to teams, and we run a series of
challenges to get the teams putting the kit together and learning how it works.

The competition is a \textbf{two day} event, typically held on a weekend in
April. After some initial tinker time, robots are pitted against each other in
an arena in a series of league matches, followed by knock-out matches to find
who has the best robot. Many robots show exciting behaviours, and most teams
significantly improve their robot code across the course of the competition
weekend.

Between the two, during the robot ``build period'', several ``tech
days'' are run by our local volunteers.  These are days where your team can
meet the other teams from the area, other SR blueshirts, and ask questions and
get advice on how to go about building their robot. The exact format and dates
of these events are locally arranged.

The latest, up to date details on all of our events can be found on the ``key
dates'' page of our
website.\footnote{\url{https://www.studentrobotics.org/key_dates}} Each event
has it's own event page, containing location, schedule and parking details.

\subsection*{Mentoring}

We believe that the best way to support teams building robots is to have a
mentor, who can give advice on anything from time management to motor and wheel
selection. Unfortunately our limited resources means we can't deliver mentoring
to everyone. We aim to send a mentor (in person) to you on a regular basis if
you can be reached by public transport from the University of Southampton or
Bristol within 45 minutes. If you're further out, we'll try to arrange a
volunteer to keep in contact with the team and provide remote mentoring,
through whatever means you find convenient.

\subsection*{Budget}

We also provide a small budget to offset the cost of buying components such as
motors and wheels. The amount of the budget changes from year to year, and an
up to date information will be on the support page of the
website\footnote{\url{https://www.studentrobotics.org/schools/support}}. We
can only reimburse purchases for which you have a receipt or invoice, and that
were for parts of the robot.

\section*{What we require from you}

As the team leader, we expect you to be a teacher, parent, or other adult who
is responsible for the competitors in the team. If you don't already have a
team of competitors, you'll need to recruit competitors to form one. We
recommend six team members, although at a minimum teams have been successful
with three members before.

The aim of Student Robotics is to educate and enthuse the competitors taking
part, and while we'd greatly appreciate you giving help and advice to the
competitors, please do not actually build their robot for them (or direct their
design). This is for their own benefit, as we expect them to make mistakes and
learn from them, as well as to make the competition a fair assessment of the
teams' ability. We trust you to decide on a fair balance of how much help to
provide.

We also rely on you to ensure that the kit is returned at the end of the
competition, or later if you arrange for a loan extension.

\subsection*{Supervision}

We are unable to provide adult supervision for competitors at any of our
events. As a result, you need to arrange adult supervision for the competitors
visiting an event. This can be yourself, or any other adult you deem suitable,
however it cannot be a competitor.

\subsection*{Health and Safety}

During the robot build period the competitors will be constructing
their robots, using whatever tools and equipment is necessary for such
construction. We expect you to be responsible for their safety during this
period. In addition, if the competitors bring any tools to the competition in
order to fix or adjust their robots there, we expect you to supervise their use
there too.

None of the equipment that makes up our electronics kit poses a risk of
electrocution, however \textbf{the batteries we supply can be a fire hazard}.
Full information on how to safely use the batteries are provided on our
documentation
pages.\footnote{\url{https://www.studentrobotics.org/docs/kit/batteries/}}
You must ensure the competitors understand that they need to treat the
batteries with respect, in particular that the charging instructions on
the website must be followed to the letter.

\subsection*{Logistics}

You will need to organise transportation to and from any of the events that we
organise, and in the case of the competition you may need to arrange overnight
accommodation and meals, too. We are not able to provide any kind of funding to
cover this, and we are also unable to reimburse the team budget against travel
costs.

If you arrange to keep the electronics kit after the competition, then you will
need to arrange for it to be couriered back to us, at your own expense. We
require registered delivery and insurance when sending kit; more information
will provided when you arrange to keep the kit.

\subsection*{Website accounts}

Many of the services on our website require an account, and you will need to
enter names and email addresses on our user management page to create accounts
for competitors. We will not share this data with third parties, although we do
export competitor names to produce certificates for attending the competition.
We only email individual competitors in the case of some irregularity with
their account (or behaviour on the website). We send infrequent bulk-emails
to brief competitors about how to get their competition ticket, what to bring
to the competition, and similar.

Please attempt to ensure that competitors do not share their accounts -- in rare
circumstances we have had to suspend competitors accounts, and in those cases
we would rather not lock the entire team out of the website.

\subsection*{Competition attendance and Media consent forms}
 
We require that everyone attending the competition fills in a media consent
form, so that we can distribute footage of the competition for publicity.
To enforce this, competitors must have a ticket to attend the competition,
which they obtain by returning a signed media consent form to you. Please
collect these forms and forward them to us (via post, or email) promptly,
and ideally well in advance of the competition.

Other attendees (such as parents, teachers, interested friends) will be able
to sign a media consent form at the door at the competition if they are
obviously over 18; everyone 18 or under will need to bring a ticket.

\end{document}
